% GENERATED FILE. Edit canonical lessons, then run npm run generate:books.
% canonical-source-hash: fnv1a64:f9e8a941d3e456f2
% canonical-lessons: TE-C35-kutumbam, TE-C35-snehitudu

\chapter{Family and Friends}
\label{ch:te-family-and-friends}
\hlchaptermodality{35}

\begin{chapteropening}
\textbf{By the end of this chapter:} \emph{I can name my family as a group in Telugu and name a friend, and say which of the two words is a Sanskrit loan and which is Telugu's own native word.}

Say \te{ఇది} \te{నా} \te{కుటుంబం} and \te{ఇతను} \te{నా} \te{స్నేహితుడు}, trace \te{కుటుంబం} and \te{స్నేహితుడు} back to Sanskrit, and set \te{స్నేహితుడు} beside the native, DEDR-attested \te{నేస్తం}.
\end{chapteropening}

\section[kuṭumbaṁ]{\te{కుటుంబం} (kuṭumbaṁ) — \textquotedblleft{}family,\textquotedblright{} the word the household words never needed}

\label{lesson:TE-C35-kutumbam}



\begin{quote}
\emph{Nāku telugu iṣṭaṁ} closed the last chapter on a sentence built from nothing but nouns. This chapter opens on the plainest nouns of all: the words for the people you would actually say that sentence to.
\end{quote}

\subsection*{You'll want to know: \te{కుటుంబం}}
\begin{quote}
\textbf{\te{కుటుంబం}} — \emph{kuṭumbaṁ} — \textbf{family}
\end{quote}

\begin{quote}
\textbf{\te{ఇది} \te{నా} \te{కుటుంబం}.} — \emph{idi nā kuṭumbaṁ.} — \textquotedblleft{}This is my family.\textquotedblright{}
\end{quote}

\textbf{\te{నా}}, \textquotedblleft{}my,\textquotedblright{} has been yours since one of the very first words you learned.

\begin{cousinweb}
\textbf{\te{కుటుంబం}} is Sanskrit \textbf{\dv{कुटुम्ब}} (\emph{kuṭumba}), \textquotedblleft{}family, household,\textquotedblright{} carried into Telugu whole — the same kind of tatsama loan as \textbf{\te{నమస్కారం}} and \textbf{\te{సంతోషం}}.

Your earlier family-words lesson gave you six words: \textbf{\te{నాన్న}}, \textbf{\te{అమ్మ}}, and the four age-graded siblings \textbf{\te{అన్న}/\te{తమ్ముడు}}, \textbf{\te{అక్క}/\te{చెల్లి}}. Every one is native Dravidian, and not one of them means \textquotedblleft{}family\textquotedblright{} — each names a single specific person. Telugu never inherited a plain word for the group itself, so when it needed one, it reached for Sanskrit. The people came first, and stayed native; the collective came later, and came borrowed — the same split you already met between \emph{tinu} and \emph{bhōjanaṁ}.
\end{cousinweb}

\subsection*{Your turn}
Say these aloud:

\begin{itemize}
\raggedright
  \item \textquotedblleft{}idi nā kuṭumbaṁ\textquotedblright{} — this is my family
  \item the six named people, then the borrowed group word — \textquotedblleft{}nānna, amma, anna, tammuḍu, akka, celli … kuṭumbaṁ\textquotedblright{}
  \item \textquotedblleft{}nāku telugu iṣṭaṁ\textquotedblright{} — once more, no verb in it
\end{itemize}

\begin{tcolorbox}[breakable,colback=teal!4,colframe=teal!35!black,title={Before you move on}]
Say \textquotedblleft{}this is my family.\textquotedblright{} (\emph{Idi nā kuṭumbaṁ}.) Is \emph{kuṭumbaṁ} native Dravidian or a loan? (\textbf{A Sanskrit loan}, like \emph{namaskāram} and \emph{santōṣam}.) Name the six family words from that earlier lesson. (\emph{Nānna, amma, anna, tammuḍu, akka, celli}.) Which of those six means \textquotedblleft{}family\textquotedblright{} as a group? (\textbf{None of them} — that's exactly the gap \emph{kuṭumbaṁ} fills.)
\end{tcolorbox}
\section[snēhituḍu]{\te{స్నేహితుడు} (snēhituḍu) — \textquotedblleft{}friend,\textquotedblright{} and the native word beside it}

\label{lesson:TE-C35-snehitudu}



\begin{quote}
\emph{Kuṭumbaṁ} named the people you're related to. This word names the people you choose.
\end{quote}

\subsection*{You'll want to know: \te{స్నేహితుడు}}
\begin{quote}
\textbf{\te{స్నేహితుడు}} — \emph{snēhituḍu} — \textbf{friend} (male; \textbf{\te{స్నేహితురాలు}}, \emph{snēhiturālu}, for a female friend)
\end{quote}

\begin{quote}
\textbf{\te{ఇతను} \te{నా} \te{స్నేహితుడు}.} — \emph{itanu nā snēhituḍu.} — \textquotedblleft{}This is my friend.\textquotedblright{}
\end{quote}

\begin{cousinweb}
\textbf{\te{స్నేహితుడు}} is built on Sanskrit \textbf{\dv{स्नेह}} (\emph{sneha}), \textquotedblleft{}affection, love\textquotedblright{} — the same kind of tatsama loan as \emph{kuṭumbaṁ} itself, plus a Telugu agentive ending.

Telugu also has a \textbf{native} word for friendship: \textbf{\te{నేస్తం}} (\emph{nēstaṁ}), catalogued in the Dravidian Etymological Dictionary alongside the related forms \emph{neyyamu}, \emph{neyyami}, and \emph{neyyari} — all clustered around friendship, love, and tenderness. \emph{Nēstaṁ} is real and still heard, but the everyday choice for \textquotedblleft{}friend\textquotedblright{} is the borrowed one, \emph{snēhituḍu} — the same pattern already found with \emph{sahāyaṁ cēyu}, where the everyday verb is built on a Sanskrit noun too.
\end{cousinweb}

\subsection*{Your turn}
Say these aloud:

\begin{itemize}
\raggedright
  \item \textquotedblleft{}itanu nā snēhituḍu\textquotedblright{} — this is my friend
  \item the borrowed word and the native one — \textquotedblleft{}snēhituḍu … nēstaṁ\textquotedblright{}
  \item family, then friend — \textquotedblleft{}kuṭumbaṁ … snēhituḍu\textquotedblright{}
\end{itemize}

\begin{tcolorbox}[breakable,colback=teal!4,colframe=teal!35!black,title={Before you move on}]
Say \textquotedblleft{}this is my friend.\textquotedblright{} (\emph{Itanu nā snēhituḍu}.) Is \emph{snēhituḍu} native or borrowed? (\textbf{A Sanskrit loan}, from \emph{sneha}, \textquotedblleft{}affection.\textquotedblright{}) What is Telugu's own native word for friendship, and is it the everyday choice? (\textbf{\te{నేస్తం}}, \emph{nēstaṁ} — real, but \textbf{not} the everyday choice; that's the borrowed word.)
\end{tcolorbox}
